\chapter{Практика}

\section{Корреляционная матрица}
Корреляционная матрица - таблица, в которой показаны зависимости между признаками.\\

Seaborn + Pandas позволяют легко построить такую матрицу:

\begin{lstlisting}
    data = pd.read_csv(data_file, sep=",")
    sns.heatmap(corr, cmap="coolwarm", center=0, square=True)
    plt.title("Correlation Matrix")
    plt.show()
\end{lstlisting}

На экране увидим:

\begin{figure}[H]
    \centering
    \includegraphics{ct.png}
    \caption{Таблица корреляции}
\end{figure}

Таблица корреляций достаточно плохая, потому что наблюдается мало зависимостей между признаками. Можно только подметить, что признак
crosses (крест) созависим с признаком saltires (наклонный крест, как на флаге Британии), а признак orange созависим с признаком animate.
Все остальные корреляции слишком малы и больше напоминают совпадение. \\


Если бы мне поставили задачу выбирать флаг, то я бы скорее всего ориентировался на то, насколько часто в среднем в флагах того или иного
региона встречается определенный цвет + ориентировался бы на фигуры. Выведем степень использования цветов/фигур по регионам:

\begin{figure}[H]
    \centering
    \includegraphics{cr.png}
    \caption{Использование цветов}
\end{figure}

\begin{figure}[H]
    \centering
    \includegraphics{cra.png}
    \caption{Использование фигур}
\end{figure}

Здесь ситуация немного лучше. Например, Африку можно легко выделить по зеленому цвету, Северную Америку по объектам живой
природы на флаге, Южную Америку отличает то, что у нее флаг разделен на 4 четверти.


\endinput